\documentclass[12pt]{article}

\usepackage[a4paper,left=20mm,right=25mm,top=25mm,bottom=25mm]{geometry}
\usepackage[utf8]{inputenc}
\usepackage[T1]{fontenc}
\usepackage[ngerman]{babel}
\usepackage{amsmath,amssymb,amsthm}
\usepackage{array}
\usepackage{booktabs}
\usepackage{tabularx}
\usepackage{hyperref}

\newtheorem{satz}{Satz}[section]
\newtheorem{definition}[satz]{Definition}
\newtheorem{lemma}[satz]{Lemma}
\newtheorem{proposition}[satz]{Proposition}
\newtheorem{korollar}[satz]{Korollar}

\newcommand{\card}[1]{\left\lvert #1\right\rvert}
\newcommand{\teamcount}[2]{\#_{#1}(#2)}

\title{Fußballspieler und Mehrheitswechsel\\
       Eine Charakterisierung durch das Medianintervall}
\author{}
\date{}

\begin{document}

\maketitle

\section{Problembeschreibung}

Es seien
\[
N=\{1,2,\ldots,n\},
\qquad
p=(p_1,p_2,\ldots,p_n)\in\{A,B\}^n.
\]
Der Eintrag $p_k$ bezeichnet die Mannschaft des Spielers an Position $k$.
Zu Beginn besitzt jede der beiden Mannschaften eine ungerade Anzahl von
Spielern.

Zwei Spieler $i,j\in N$ mit $i<j$ und $p_i=p_j$ möchten gemeinsam zur
anderen Mannschaft wechseln. Die Spieler an den Positionen $i$ und $j$
stimmen nicht ab. Jeder andere Spieler stimmt genau dann mit \emph{Ja}, wenn
seine Position strikt zwischen $i$ und $j$ liegt; andernfalls stimmt er mit
\emph{Nein}.

Für $c\in\{A,B\}$ definieren wir daher
\begin{align*}
Y_c(i,j;p)
  &=\#\{k\in N\mid i<k<j,\ p_k=c\},\\
Z_c(i,j;p)
  &=\#\{k\in N\mid k<i\text{ oder }j<k,\ p_k=c\}.
\end{align*}

\begin{definition}[Zulässiger Wechsel]
Ein Paar $(i,j)$ heißt in $p$ \emph{zulässig}, wenn
\[
i<j,\qquad p_i=p_j
\]
und beide Mannschaften dem Wechsel mit strikter Mehrheit zustimmen, also
\[
Y_A(i,j;p)>Z_A(i,j;p),
\qquad
Y_B(i,j;p)>Z_B(i,j;p).
\]
Für ein zulässiges Paar bezeichnet $\tau_{i,j}(p)$ den Vektor, der aus $p$
durch Vertauschen der Mannschaftszugehörigkeit an den Positionen $i$ und $j$
entsteht:
\[
(\tau_{i,j}(p))_k=
\begin{cases}
B,&k\in\{i,j\}\text{ und }p_k=A,\\
A,&k\in\{i,j\}\text{ und }p_k=B,\\
p_k,&k\notin\{i,j\}.
\end{cases}
\]
\end{definition}

\begin{definition}[Erreichbarkeit]
Für $p,q\in\{A,B\}^n$ schreiben wir $p\sim q$, wenn $q$ aus $p$ durch eine
endliche Folge zulässiger Wechsel entsteht. Die leere Folge ist erlaubt, also
gilt stets $p\sim p$.
\end{definition}

Gesucht ist ein Algorithmus, der für ein gegebenes Anfangsmuster $p$ und ein
Zielmuster $q$ entscheidet, ob $p\sim q$ gilt.

\section{Grundlegende Beobachtungen}

\begin{proposition}[Parität der Mannschaftsgrößen]
Bei jedem Wechsel ändert sich die Größe jeder Mannschaft um genau zwei.
Insbesondere bleiben die Paritäten beider Mannschaftsgrößen invariant.
\end{proposition}

\begin{proof}
Die beiden wechselnden Spieler gehören vor dem Wechsel derselben Mannschaft
an. Diese verliert zwei Spieler, während die andere Mannschaft zwei Spieler
gewinnt.
\end{proof}

Da beide Mannschaftsgrößen in $p$ ungerade sind, ist $n$ die Summe zweier
ungerader Zahlen und somit gerade. Außerdem kann ein Zielmuster mit einer
geraden Mannschaftsgröße niemals erreicht werden.

\begin{proposition}[Umkehrbarkeit]
Jeder zulässige Wechsel ist umkehrbar.
\end{proposition}

\begin{proof}
Nach dem Wechsel sind dieselben Spieler $i$ und $j$ erneut in derselben
Mannschaft. Alle abstimmenden Spieler und ihre Positionen sind unverändert;
folglich bleiben auch sämtliche Ja- und Nein-Stimmen unverändert. Der
Rückwechsel wird daher von beiden Mannschaften mit derselben Mehrheit
angenommen.
\end{proof}

\section{Medianen und Medianintervall}

\begin{definition}[Mannschaftsmedian]
Für ein Muster $s\in\{A,B\}^n$ sei
\[
P_c(s)=\{k\in N\mid s_k=c\}
\qquad(c\in\{A,B\}).
\]
Ist $\card{P_c(s)}=2h_c+1$, so heißt die Position des $(h_c+1)$-ten
Vorkommens von $c$ der \emph{Median} der Mannschaft $c$ und wird mit
$m_c(s)$ bezeichnet.
\end{definition}

\begin{definition}[Mediansignatur]
Besitzen beide Mannschaften in $s$ ungerade Größe, so setzen wir
\begin{align*}
L(s)&=\min\{m_A(s),m_B(s)\},\\
R(s)&=\max\{m_A(s),m_B(s)\}
\end{align*}
und definieren die \emph{Mediansignatur}
\[
\Sigma(s)=
\bigl(L(s),R(s),s_{L(s)}s_{L(s)+1}\cdots s_{R(s)}\bigr).
\]
Das Teilwort $s_{L(s)}\cdots s_{R(s)}$ heißt \emph{Medianintervall}.
\end{definition}

\begin{lemma}[Beide Medianen liegen im Wechselintervall]
\label{lem:medianen-im-intervall}
Ist $(i,j)$ in $s$ zulässig, so gilt
\[
i<m_A(s)<j
\qquad\text{und}\qquad
i<m_B(s)<j.
\]
\end{lemma}

\begin{proof}
Zunächst betrachten wir die Mannschaft $c=s_i=s_j$ der beiden Endpunkte.
Sie besitze $2h+1$ Spieler. Die Endpunkte seien das $r$-te und das $t$-te
Vorkommen von $c$, wobei $r<t$ gilt. Dann erhält diese Mannschaft
\[
t-r-1
\]
Ja-Stimmen und
\[
(r-1)+(2h+1-t)
\]
Nein-Stimmen. Aus der notwendigen Ungleichung
\[
t-r-1>(r-1)+(2h+1-t)
\]
folgt $t-r\ge h+1$. Daher kann weder $r\ge h+1$ noch $t\le h+1$ gelten;
also liegt das $(h+1)$-te Vorkommen strikt zwischen den Endpunkten.

Sei nun $d\ne c$ die andere Mannschaft mit $2k+1$ Spielern. Alle ihre
Spieler stimmen ab. Eine strikte Ja-Mehrheit erfordert mindestens $k+1$
Vorkommen von $d$ im offenen Intervall $(i,j)$. Die darin enthaltenen
Vorkommen bilden einen zusammenhängenden Block in der geordneten Liste aller
$d$-Vorkommen. Jeder solche Block der Länge mindestens $k+1$ enthält das
$(k+1)$-te Vorkommen. Somit liegt auch $m_d(s)$ strikt zwischen $i$ und $j$.
\end{proof}

\begin{lemma}[Invarianz der Mediansignatur]
\label{lem:signatur-invariant}
Ist $s'=\tau_{i,j}(s)$ ein zulässiger Wechsel, so gilt
\[
\Sigma(s')=\Sigma(s).
\]
\end{lemma}

\begin{proof}
Nach Lemma~\ref{lem:medianen-im-intervall} liegen beide Medianen strikt
zwischen $i$ und $j$. Die abgebende Mannschaft verliert somit einen Spieler
links und einen Spieler rechts von ihrem Median. Die aufnehmende Mannschaft
gewinnt entsprechend einen Spieler auf jeder Seite ihres Medians. In beiden
Mannschaften besitzt der bisherige Median danach weiterhin gleich viele
Mitspieler links und rechts von sich. Beide Medianpositionen bleiben also
unverändert.

Ferner gilt
\[
i<L(s)\le R(s)<j.
\]
Innerhalb von $[L(s),R(s)]$ wird daher kein Eintrag verändert. Sowohl die
Grenzen als auch das vollständige Teilwort des Medianintervalls bleiben
gleich.
\end{proof}

Damit ist bereits die notwendige Richtung bewiesen:
\begin{equation}
p\sim q\quad\Longrightarrow\quad\Sigma(p)=\Sigma(q).
\label{eq:notwendigkeit}
\end{equation}

\section{Konstruktion außerhalb des Medianintervalls}

\begin{lemma}[Randgestützte Wechsel]
\label{lem:randwechsel}
Es sei $[L,R]$ das Medianintervall eines Musters $s$.
\begin{enumerate}
\item Gilt $i<L$ und $s_i=s_n$, so ist $(i,n)$ zulässig.
\item Gilt $j>R$ und $s_1=s_j$, so ist $(1,j)$ zulässig.
\end{enumerate}
\end{lemma}

\begin{proof}
Wegen der Symmetrie genügt es, die erste Behauptung zu zeigen. Die gemeinsame
Mannschaft der Endpunkte besitze $2h+1$ Spieler. Da $i$ strikt vor ihrem
Median liegt, befinden sich höchstens $h-1$ weitere Mitglieder dieser
Mannschaft vor $i$. Diese stimmen mit Nein. Von den übrigen $2h-1$ Wählern
dieser Mannschaft stimmen daher mindestens $h$ mit Ja. Es gibt also strikt
mehr Ja- als Nein-Stimmen.

Die andere Mannschaft besitze $2k+1$ Spieler. Da $i$ auch vor deren Median
liegt, befinden sich höchstens $k$ ihrer Spieler vor $i$. Keiner ihrer Spieler
ist ein Endpunkt. Somit stimmen höchstens $k$ mit Nein und mindestens $k+1$
mit Ja. Auch diese Mannschaft stimmt dem Wechsel zu.
\end{proof}

\begin{lemma}[Randkonstruktion]
\label{lem:randkonstruktion}
Seien $i<L$ und $j>R$ mit $s_i=s_j$. Dann existiert eine Folge von höchstens
drei zulässigen Wechseln, deren Gesamteffekt genau die Einträge an den
Positionen $i$ und $j$ umschaltet. Alle anderen Einträge werden wieder auf
ihren ursprünglichen Wert zurückgesetzt.
\end{lemma}

\begin{proof}
Falls $i=1$ oder $j=n$ gilt, ist der direkte Wechsel nach
Lemma~\ref{lem:randwechsel} zulässig. Es sei daher
\[
1<i<L\le R<j<n
\]
und $x=s_i=s_j$. Je nach Belegung der Randpositionen führen wir die folgende
Folge aus; die Wechsel sind von links nach rechts zu lesen.

\begin{center}
\begin{tabularx}{\textwidth}{@{}>{\raggedright\arraybackslash}X
                                  >{\raggedright\arraybackslash}X@{}}
\toprule
Randwerte & Folge der Wechsel \\
\midrule
$s_1=s_n=x$
  & $\tau_{i,n},\ \tau_{1,j},\ \tau_{1,n}$ \\
$s_1=s_n\ne x$
  & $\tau_{1,n},\ \tau_{i,n},\ \tau_{1,j}$ \\
$s_1\ne s_n$ und $s_n=x$
  & $\tau_{i,n},\ \tau_{1,n},\ \tau_{1,j}$ \\
$s_1\ne s_n$ und $s_1=x$
  & $\tau_{1,j},\ \tau_{1,n},\ \tau_{i,n}$ \\
\bottomrule
\end{tabularx}
\end{center}

Bei jedem einzelnen Schritt sind die beiden angegebenen Endpunkte gleich.
Nach Lemma~\ref{lem:randwechsel} ist der Schritt somit zulässig. Nach
Lemma~\ref{lem:signatur-invariant} bleibt das Medianintervall auch in allen
Zwischenzuständen dasselbe. Eine direkte Kontrolle der vier Fälle zeigt, dass
die Randpositionen $1$ und $n$ am Ende wieder ihren Anfangswert besitzen,
während genau $i$ und $j$ umgeschaltet wurden.
\end{proof}

\begin{lemma}[Binäres Bilanzlemma]
\label{lem:bilanz}
Es seien $I$ und $J$ zwei nichtleere, disjunkte Mengen binärer Positionen.
Als abstrakte Operation dürfe man einen Eintrag aus $I$ und einen Eintrag aus
$J$ genau dann gleichzeitig umschalten, wenn ihre aktuellen Werte gleich
sind. Zwei Belegungen $u$ und $v$ sind genau dann durch solche Operationen
verbunden, wenn
\[
\teamcount{A}{u|_I}-\teamcount{A}{u|_J}
=
\teamcount{A}{v|_I}-\teamcount{A}{v|_J}
\]
gilt.
\end{lemma}

\begin{proof}
Jede Operation erhöht beide Anzahlen von A um eins oder verringert beide um
eins. Die Differenz ist daher invariant.

Für die umgekehrte Richtung vergleichen wir eine aktuelle Belegung mit $v$.
Gibt es auf verschiedenen Seiten zwei Fehlstellen mit gleichem aktuellem
Wert, so schalten wir diese gemeinsam um und beseitigen damit beide
Fehlstellen.

Gibt es auf einer Seite zwei Fehlstellen mit verschiedenen aktuellen Werten,
so wählen wir eine beliebige Position der anderen Seite als Zwischenspeicher.
Zuerst schalten wir diejenige Fehlstelle um, deren Wert dem Speicherwert
entspricht. Danach stimmt der umgeschaltete Speicherwert mit der zweiten
Fehlstelle überein, sodass wir beide gemeinsam umschalten können. Beide
Fehlstellen sind nun korrigiert, und der Speicher besitzt wieder seinen
ursprünglichen Wert.

Angenommen, keine dieser beiden Reduktionen wäre trotz vorhandener
Fehlstellen möglich. Dann hätten alle Fehlstellen jeder Seite jeweils
denselben aktuellen Wert, und die Werte der beiden Seiten wären verschieden.
Die Beiträge der beiden Seiten zur behaupteten Bilanzgleichung hätten dann
entgegengesetzte, von null verschiedene Vorzeichen. Dies widerspricht der
Bilanzgleichung. Folglich lässt sich eine der beiden Reduktionen so lange
anwenden, bis keine Fehlstelle mehr vorhanden ist.
\end{proof}

\section{Hauptsatz und Algorithmus}

\begin{satz}[Vollständige Charakterisierung]
\label{satz:charakterisierung}
Es seien $p,q\in\{A,B\}^n$, und beide Mannschaften mögen sowohl in $p$ als
auch in $q$ ungerade Größe besitzen. Dann gilt
\[
p\sim q
\quad\Longleftrightarrow\quad
\Sigma(p)=\Sigma(q).
\]
\end{satz}

\begin{proof}
Die notwendige Richtung folgt aus
Lemma~\ref{lem:signatur-invariant} beziehungsweise aus
Gleichung~\eqref{eq:notwendigkeit}.

Für die hinreichende Richtung gelte $\Sigma(p)=\Sigma(q)$. Das gemeinsame
Medianintervall sei $[L,R]$ und das gemeinsame Teilwort sei $C$. Wir setzen
\[
I=\{1,\ldots,L-1\},
\qquad
J=\{R+1,\ldots,n\}.
\]

Zunächst zeigen wir die Bilanzgleichung aus Lemma~\ref{lem:bilanz}. Liegt der
A-Median bei $L$, so befinden sich in jedem Muster $s$ mit dieser Signatur
gleich viele A-Spieler links und rechts von Position $L$. Daher gilt
\[
\teamcount{A}{s|_I}
=
\bigl(\teamcount{A}{C}-1\bigr)+\teamcount{A}{s|_J}
\]
und somit
\[
\teamcount{A}{s|_I}-\teamcount{A}{s|_J}
=\teamcount{A}{C}-1.
\]
Liegt der A-Median dagegen bei $R$, erhält man analog
\[
\teamcount{A}{s|_I}-\teamcount{A}{s|_J}
=1-\teamcount{A}{C}.
\]
In beiden Fällen hängt die Differenz ausschließlich vom gemeinsamen Teilwort
$C$ ab. Sie besitzt folglich für $p$ und $q$ denselben Wert.

Sind $I$ und $J$ nichtleer, so liefert Lemma~\ref{lem:bilanz} eine Folge
abstrakter Umschaltungen gleicher Werte, welche die Einträge außerhalb von
$C$ von $p$ nach $q$ überführt. Jede abstrakte Umschaltung lässt sich nach
Lemma~\ref{lem:randkonstruktion} durch höchstens drei zulässige Wechsel
realisieren, ohne das Medianintervall zu verändern. Somit gilt $p\sim q$.

Es bleibe der Randfall $I=\varnothing$. Die Mannschaft, deren Median an
Position $1$ liegt, besitzt kein Mitglied vor ihrem Median. Da ein Median
gleich viele Mannschaftsmitglieder links und rechts von sich hat, besteht
diese Mannschaft aus genau einem Spieler. Sie kann daher in $J$ nicht mehr
vorkommen; sämtliche Einträge in $J$ sind eindeutig als die andere
Mannschaft bestimmt. Aus der gemeinsamen Signatur folgt bereits $p=q$. Der
Fall $J=\varnothing$ ist symmetrisch.
\end{proof}

\begin{korollar}[Entscheidungskriterium für die Eingabe]
Für ein gültiges Anfangsmuster $p$ gilt genau dann $p\sim q$, wenn beide
Mannschaften auch in $q$ ungerade Größe besitzen und
\[
\Sigma(p)=\Sigma(q)
\]
gilt.
\end{korollar}

\begin{proof}
Die Notwendigkeit der ungeraden Mannschaftsgrößen folgt aus der
Paritätsinvarianz. Sind beide Größen ungerade, ist
Satz~\ref{satz:charakterisierung} anwendbar.
\end{proof}

\subsection*{Algorithmus}

Für jedes der beiden Muster wird die Mediansignatur in einem linearen
Durchlauf bestimmt:
\begin{enumerate}
\item Speichere die Positionen aller A-Spieler und aller B-Spieler.
\item Wähle aus jeder Liste das mittlere Element als Mannschaftsmedian.
\item Setze $L$ und $R$ auf die kleinere beziehungsweise größere
      Medianposition.
\item Gib das Tupel $(L,R,s_Ls_{L+1}\cdots s_R)$ als Signatur zurück.
\end{enumerate}

Besitzt in $q$ eine Mannschaft gerade Größe, lautet die Antwort sofort
\texttt{NO}. Andernfalls lautet sie genau dann \texttt{YES}, wenn die beiden
Signaturen übereinstimmen.

\begin{korollar}[Komplexität]
Die Entscheidung benötigt $O(n)$ Zeit und in der beschriebenen
Implementierung $O(n)$ zusätzlichen Speicher. Über alle Testfälle beträgt die
Laufzeit $O(\sum n)$.
\end{korollar}

\begin{proof}
Jedes Zeichen wird nur konstant oft betrachtet. Die beiden Positionslisten
enthalten zusammen genau $n$ Einträge. Alle weiteren Operationen sind
konstant oder proportional zur Länge des Medianintervalls, also ebenfalls
$O(n)$.
\end{proof}

\end{document}
